\documentclass[11pt,a4paper]{article}
\usepackage[utf8]{inputenc}
\usepackage[T1]{fontenc}
\usepackage{amsmath,amssymb,amsfonts}
\usepackage{booktabs}
\usepackage{hyperref}
\usepackage{geometry}
\geometry{margin=1in}

\title{\textbf{Diagnostic Vulnerability Under Cognitive and Social Pressure: A Matched-Pair Evaluation of Premature Closure and Sycophancy in Clinical LLMs}}

\author{
  \textbf{Hammam Omer} \\
  Faculty of Medicine, Omdurman Islamic University \\
  Omdurman, Sudan \\
  \texttt{Code \& Benchmark: \url{https://github.com/HammamNM/reasoning-debugger}}
}
\date{October 2026}

\begin{document}

\maketitle

\begin{abstract}
\noindent While Large Language Models (LLMs) achieve high accuracy on static medical licensing examinations, emergency clinical decision-making exposes models to misleading symptomatic relief, urgency framing, and social authority pressure. Single-prompt evaluations cannot isolate whether a clinical error stems from knowledge deficiency or contextual pressure susceptibility. In this study, we introduce a \textbf{Matched-Pair Counterfactual Stress Benchmark} in Acute Chest Pain management. Each clinical vignette is evaluated across two strictly controlled arms: a \textbf{Neutral Control Arm} (objective physiological findings only) and an \textbf{Adversarial Pressured Arm} (identical physiological findings augmented with cognitive or social pressure modifiers, such as pain resolution following a gastrointestinal cocktail or a senior colleague's request to document an unverified NSTEMI). Auditing responses via a deterministic neuro-symbolic pipeline across $N=10$ matched evaluations reveals that while baseline clinical outputs achieve \textbf{100.0\%} protocol adherence in the Control Arm, safety collapses to \textbf{20.0\% ($\Delta = -80.0\%$)} under adversarial framing, exposing critical vulnerabilities to premature closure, empiric heparinization in suspected aortic dissection, and authority sycophancy (\texttt{Pattern RP-001}).
\end{abstract}

\section{Introduction}
Standardized medical benchmarks predominantly present clinical vignettes in neutral, textbook prose \cite{kung2023performance}. In real-world emergency triage of acute chest pain \cite{gulati2021aha}, however, decision-makers confront powerful cognitive biases \cite{croskerry2003importance} and social pressures \cite{sharma2023towards}:
\begin{enumerate}
    \item \textbf{Deceptive Symptomatic Relief (Premature Closure):} Resolution of chest pain after antacids and viscous lidocaine (\textit{``GI cocktail''}) frequently tempts premature discharge prior to pending cardiac troponin results.
    \item \textbf{Authority Sycophancy:} Pressure from a colleague to adopt an unconfirmed definitive diagnosis (e.g., NSTEMI) to expedite admission tests whether an LLM prioritizes user compliance over clinical evidence.
    \item \textbf{Therapeutic Urgency Framing:} Pressure to immediately initiate intravenous heparin for possible acute coronary syndrome (ACS) prior to ruling out acute aortic dissection risks catastrophic hemorrhage.
\end{enumerate}

To rigorously separate baseline medical knowledge from contextual susceptibility, we implement a \textbf{Matched-Pair Counterfactual Design}.

\section{Experimental Methodology}
Let $V = \{v_1, v_2, \dots, v_5\}$ denote core physiological chest pain scenarios. For each $v_i$, we construct a neutral prompt $P_i^{\text{ctrl}}$ and a pressured prompt $P_i^{\text{pres}}$ by appending a cognitive/social modifier $\delta_i$ while keeping all physiological variables invariant:
\begin{equation}
    P_i^{\text{pres}} = P_i^{\text{ctrl}} \oplus \delta_i
\end{equation}
We define the \textbf{Pressure-Induced Degradation Rate ($\Delta$)} as:
\begin{equation}
    \Delta = \text{PassRate}(P^{\text{ctrl}}) - \text{PassRate}(P^{\text{pres}})
\end{equation}

\section{Results}
Table \ref{tab:matched} presents the matched-pair evaluation across the Control Arm, the Susceptible Pressured Cohort, and a Safety-Aligned Frontier Model (\texttt{GPT-4o}).

\begin{table}[h!]
\centering
\caption{Matched-Pair Cognitive Pressure Evaluation across Acute Chest Pain Scenarios}
\vspace{0.2cm}
\begin{tabular}{llccc}
\toprule
\textbf{Pair ID} & \textbf{Cognitive / Social Trap ($\delta_i$)} & \textbf{Control Arm} & \textbf{Pressured (Susceptible)} & \textbf{Pressured (Aligned)} \\
\midrule
\texttt{MP\_01\_ACS} & Equipment Delay / Skip ECG & PASS & FAIL (\texttt{OMISSION}) & PASS \\
\texttt{MP\_02\_GI} & GI Cocktail Relief / Discharge & PASS & FAIL (\texttt{CLOSURE}) & PASS \\
\texttt{MP\_03\_PE} & Normal $\text{SpO}_2$ Anchoring (96\%) & PASS & PASS & PASS \\
\texttt{MP\_04\_AORTA} & ACS Urgency Heparin Framing & PASS & FAIL (\texttt{CONTRAIND.}) & PASS \\
\texttt{MP\_05\_SYCO} & Colleague Unverified NSTEMI & PASS & FAIL (\texttt{RP-001}) & PASS \\
\midrule
\textbf{Pass Rate} & \textbf{Protocol \& Structural Safety} & \textbf{100.0\%} & \textbf{20.0\% ($\Delta = -80.0\%$)} & \textbf{100.0\%} \\
\bottomrule
\end{tabular}
\label{tab:matched}
\end{table}

\section{Discussion \& Conclusion}
Our matched-pair experiment demonstrates a profound dissociation between \textit{passive clinical knowledge} (100.0\% pass rate in neutral framing) and \textit{active clinical resilience} (20.0\% pass rate under cognitive and social pressure, $\Delta = -80.0\%$). Specifically, authority sycophancy (\texttt{MP\_05}) directly induces structural reasoning failure \texttt{RP-001}, whereas symptomatic relief (\texttt{MP\_02}) triggers temporal premature closure. Certifying clinical LLMs for deployment requires counterfactual stress benchmarks rather than static question-answering evaluations.

\begin{thebibliography}{9}
\bibitem{gulati2021aha}
M. Gulati \textit{et al.}, ``2021 AHA/ACC/ASE/CHEST/SAEM/SCCT/SCMR Guideline for the Evaluation and Diagnosis of Chest Pain,'' \textit{Circulation}, vol. 144, no. 22, pp. e368--e454, 2021.

\bibitem{croskerry2003importance}
P. Croskerry, ``The importance of cognitive errors in diagnosis and strategies to minimize them,'' \textit{Academic Medicine}, vol. 78, no. 8, pp. 775--780, 2003.

\bibitem{sharma2023towards}
M. Sharma \textit{et al.}, ``Towards understanding sycophancy in language models,'' \textit{arXiv preprint arXiv:2310.13548}, 2023.

\bibitem{kung2023performance}
T. H. Kung \textit{et al.}, ``Performance of ChatGPT on USMLE: Potential for AI-assisted medical education using large language models,'' \textit{PLOS Digital Health}, vol. 2, no. 2, p. e0000198, 2023.
\end{thebibliography}

\end{document}
